\documentclass[sigconf,authorversion,nonacm]{acmart}

%%% ADDED PACKAGES %%%
\usepackage{url}
\usepackage[belowskip=0pt,aboveskip=5pt]{subcaption}
\usepackage{listings}
\usepackage{enumitem}
\usepackage{booktabs}
\usepackage{xspace}
\usepackage{makecell}
\usepackage{amsmath}
\usepackage{pifont}
\usepackage{balance}
\usepackage{soul}
\usepackage{microtype}    
\usepackage[ruled,vlined,linesnumbered]{algorithm2e}
\usepackage[noend]{algpseudocode}
\usepackage{amsthm}
\usepackage{comment}
\usepackage{xr}
\usepackage{color}
\usepackage{setspace}
\usepackage{listings}
\usepackage{lipsum}
\usepackage{adjustbox}
\usepackage{array,multirow,graphicx}
\usepackage{float}
\usepackage{graphicx}
\usepackage{algorithm}
\usepackage{amsfonts}
% \documentclass{article}
\usepackage{tikz}
\usetikzlibrary{shapes, arrows.meta, positioning, fit, backgrounds, calc}

\lstdefinestyle{code}{
  backgroundcolor=\color{gray!4},
  commentstyle=\color{codegray},
  keywordstyle=\color{codepurple},
  numberstyle=\tiny\color{codegray},
  stringstyle=\color{codegray},
  basicstyle=\ttfamily\footnotesize,
  breakatwhitespace=false,         
  breaklines=true,                 
  captionpos=b,                    
  keepspaces=true,                 
  numbers=left,                    
  numbersep=5pt,                  
  showspaces=false,                
  showstringspaces=false,
  showtabs=false,                  
  tabsize=2,
  % escapeinside=||
}
\lstset{style=code}

\lstset{escapeinside={<@}{@>}}
%%% ADDED Commands %%%
\definecolor{comment}{rgb}{0,0.6,0}
\definecolor{dgray}{gray}{0.25}
\definecolor{strings}{RGB}{0,0,128}
\definecolor{backcolour}{rgb}{0.95,0.95,0.92}
\definecolor{keywords}{RGB}{127,0,85}
\definecolor{darkred}{RGB}{139,0,0}
\definecolor{darkyellow}{RGB}{204,204,0}
\definecolor{darkblue}{rgb}{0.0, 0.0, 0.55}
\definecolor{vividviolet}{rgb}{0.62, 0.0, 1.0}
\definecolor{fuchsia}{rgb}{1.0, 0.0, 1.0}
\definecolor{shockingpink}{rgb}{0.99, 0.06, 0.75}
\definecolor{darkBlue}{RGB}{0, 0, 205}
\lstset{breaklines=true}
\captionsetup[lstlisting]{font={tiny}}
\definecolor{problemblue}{RGB}{100,134,158}
\definecolor{idiomsgreen}{RGB}{0,162,0}
\definecolor{exercisebgblue}{rgb}{0,  .69,  .941}
\definecolor{deepgreen}{rgb}{0.0, 0.5, 0.0}
\definecolor{codegreen}{rgb}{0,0.6,0}
\definecolor{codegray}{rgb}{0.5,0.5,0.5}
\definecolor{codepurple}{rgb}{0.58,0,0.82}
\definecolor{backcolour}{rgb}{0.95,0.95,0.92}
\definecolor{redColor}{RGB}{255,0,0}
\definecolor{Gray}{gray}{0.1}
\definecolor{CadetBlue}{RGB}{243.0,42.1,134.0}

%%% Packages for Algorithms %%%
\usepackage[linesnumbered,ruled,vlined]{algorithm2e}
\SetKwInput{KwInput}{Input}                % Set the Input
\SetKwInput{KwInputs}{Inputs}                % Set the Input
\SetKwInput{KwOutput}{Output}              % set the Output

\begin{document}

\title{CS 527 Research Project Report: Group-18}

\author{Ang Li, Do Hoon Kim}
\affiliation{
    \institution{University of Illinois at Urbana-Champaign}
}
\email{[angl5,dhk8]@illinois.edu}

\begin{abstract}
    % [proposal] One or more opening sentences on the importance of problem. 
    % One of more statements about limitations of existing approaches. 
    % One or more statements on how your proposed technique bridges the gap.
    % \newline
    % To enhance the efficiency of long-term test sets, we address complex code dependencies overlooked by existing techniques through GNN learning.
    !!! Fix later !!!
    In large-scale industrial software development, long-running test suites (LRTS) often take 6.5 hours, causing severe bottlenecks in continuous integration (CI) pipeline. To address this, history-based prioritization techniques like LatestFail+CC have been adopted due to low overhead. However, these methods inherently struggle to detect 'First Failures' because they rely on past execution data and ignore the complex dependencies between code and tests. To bridge this gap, we propose GNN-TCP, a new test prioritization method based on Relational Graph Convolution Networks (R-GCN). By constructing a Code Dependency Graph (CDG) that integrates method calls and test coverage, our model learns the structural propagation of code change impacts to predict failure probabilities. We evaluate GNN-TCP on LRTS dataset comprising 10 large-scale open-source projects. This evaluation aims to demonstrate that the dependency-aware approach significantly outperforms 59 baseline techniques in terms of APFDc. By reducing the training and inference overhead compared to RL or mutation-based SOTA models, GNN-TCP proves to be a more practical solution for large-scale industrial CI environment.
\end{abstract}

\maketitle

\section{Artifact Information}

Link to the GitHub Repository: \url{[https://github.com/Dohoon-Dean-Kim/CS527-group18]} 
    % Waiting for the professor's response

%LRTS dataset https://zenodo.org/records/12662090 | https://github.com/lrtsuser/LRTS
\raggedbottom
\section{Introduction}

% [proposal] One paragraph about the importance of the problem

% [proposal] One or two paragraphs about state-of-the-art techniques that have tried to solve the problem

% [proposal] One paragraph about the research gap (why and how existing techniques fail to properly treat the problem)

% [proposal] One paragraph on how the proposed technique will bridge the research gap

% [proposal] One paragraph about the experimental plan

% One paragraph about the results of experiments

% A list of contributions your technique/paper makes.

% Industrial projects frequently encounter Long-Running Test Suites (LRTS) that take over 6.5 hours to execute. Therefore, efficient test case prioritization is essential to provide timely feedback by running critical test first.

% For large-scale systems like those in the LRTS dataset, traditional methods are often too slow because of high computational overhead. Therefore, recent study focuses on history-based techniques which are faster and lighter. According to the ISSTA '24, the current state-of-the-art is LatestFail+CC. This hybrid approach combines recent failure history with a lightweight coverage metric, offering the best balance between speed and fault detection. It serves as the primary baseline for our research. In addition, recent academic papers introduced high performance models such as XGRL-TCP[2], which use reinforcement learning, and mutation-aware GNNs [3].

% The latest technique, LatestFail+CC, is simple but has the limitation of relying on past history. It assumes that tests which failed recently are likely to fail again. As a result, it may ignore the influence of new modifications in stable areas. Therefore, existing techniques are vulnerable to detecting 'First Failures' because they fail to capture the complex dependencies between tests and code. Because they do not consider how changes spread through the codebase, these methods fail to identify risks with no prior history. Due to the high computational costs arising from repetitive reinforcement learning training or mutation analysis, it is impractical for a 6.5 hour LRTS cycle.

% We perform sophisticated dependency-based failure probability prediction by learning the propagation paths of code change impacts through Graph Neural Networks (GNNs). By utilizing a Relational Graph Convolutional Network, our model aggregates risk features from modified methods and propagates them to potential test nodes, assigning higher failure probabilities to tests affected by the changes. This approach aims to achieve a better balance between the speed of history-based methods and the accuracy of complex GNN models.

% We will quantitatively validate the effectiveness of our GNN model using Long-Running Test Suites (LRTS) dataset, which includes 10 large-scale projects. We compare our approach against 59 baseline techniques, measuring performance using Average Percentage of Faults Detected (APFD) and its cost-cognizant variant (APFDc). Our evaluation aims to demonstrate improved efficiency over traditional baselines and recent SOTA models [2, 3], particularly for detecting first-time failures.

% % A list of contributions your technique/paper makes.
% A list of contributions is as follows:
% \begin{itemize}
%     \item We apply the Relational Graph Convolutional Network (R-GCN) for the first time to test prioritization. This model learns how code changes propagate through the graph to predict failure probabilities.
%     \item We overcome the limitations of history-based techniques by detecting ‘First Failures’. Unlike existing methods relying on past data, our approach identifies risks based on actual code dependencies.
%     \item We validated our approach using the LRTS dataset (10 large-scale projects). Benchmarking against 59 baseline techniques and recnet GNN-based SOTAs demonstrated our model's robustness and effectiveness.
% \end{itemize}

%%% Introduction changed

Industrial projects frequently encounter Long-Running Test Suites (LRTS) that take over 6.5 hours to execute. In such environments, identifying failing builds early in the continuous integration (CI) pipeline is essential to save computational resources and provide timely feedback to developers. For large-scale systems like those in the LRTS dataset, traditional dynamic analysis methods are often too slow because of high computational overhead. Consequently, recent studies have heavily focused on history-based techniques or lightweight coverage metrics to balance speed and fault detection.

To further optimize this process, our initial intuition was to perform sophisticated Test Case Prioritization (TCP) by learning the propagation paths of code change impacts through a Relational Graph Convolutional Network (R-GCN). We aimed to construct a complete Code Dependency Graph (CDG) utilizing method-level dynamic coverage to precisely target 'First Failures'. However, during the empirical data extraction phase on the LRTS dataset comprising 10 large-scale open-source projects, we encountered severe technical constraints. Specifically, accessing legacy build artifacts and successfully injecting dynamic profiling agents (e.g., JaCoCo) caused frequent Java/Scala version conflicts and build failures. Securing consistent method-level dynamic data across the entire suite proved to be highly impractical and unscalable in a real-world industrial context.

Given these constraints, this paper pivots from dynamic coverage-based test prioritization to an empirical evaluation of \textbf{Build Failure Prediction} using static analysis and CI metadata. Rather than relying on incomplete or unstable dynamic execution traces, we examine a lightweight Graph Convolutional Network (GCN) encoder that operates on static source code call graphs. By combining structural node features with per-build temporal metadata, we assess whether this approach can predict failing builds ($num\_fail\_class > 0$) as a preliminary gatekeeper.

Our experimental results indicate that while the GCN model provides a computationally efficient framework, it achieves modest performance, particularly exhibiting a low recall. These underwhelming results suggest that static features alone are insufficient to capture the dynamic execution risks of complex industrial projects. These findings serve as a valuable empirical baseline, highlighting the significant performance gap between static and dynamic analysis and confirming the necessity of high-fidelity dynamic coverage for reliable build failure prediction in large-scale environments.

Our main contributions are as follows:
\begin{itemize}
    \item We provide a detailed empirical analysis of the severe technical constraints and data sparsity encountered when attempting to extract method-level dynamic dependencies from large-scale industrial projects.
    \item We implement a GCN-style encoder combining static call graphs with CI metadata and demonstrate the limitations of a purely static approach in build failure prediction.
    \item We analyze the performance gap between static and dynamic features, offering insights into the necessity of dynamic analysis for robust CI integration in legacy LRTS environments.
\end{itemize}


% One paragraph about the results of experiments

% A list of contributions your technique/paper makes. 
\section{Proposed Technique}

\subsection{Approach Overview}
% [proposal] Describe different components and stages of your technique using text/figure/example/flowchart/etc. 
Our initial design was a test-level GNN-TCP system: construct a graph connecting changed methods to covered tests, propagate change impact through the graph, and rank tests by predicted failure probability. In practice, robust method-level dynamic coverage was not available for the full LRTS setting. The implemented technique is therefore a static build-failure baseline. It has three stages: (1) construct one static call graph per project, (2) encode each project graph using a two-layer Graph Convolutional Network (GCN), and (3) concatenate the resulting project embedding with per-build CI metadata to predict whether a build contains at least one failing test class.

\subsection{Static Graph Construction}

For each project $p$, we represent the source code as a static directed call graph
\[
G_p=(V_p,E_p),
\]
where each node $v\in V_p$ is a method and each directed edge $(u,v)\in E_p$ indicates that method $u$ invokes method $v$. The graph is extracted from Java Abstract Syntax Trees by identifying \texttt{MethodDeclaration} and \texttt{MethodInvocation} nodes.

Let $d^-_v$ and $d^+_v$ denote the in-degree and out-degree of node $v$ in the directed call graph. The implemented node feature vector is purely structural:
\[
x_v =
\begin{bmatrix}
\log(1+d^-_v) &
\log(1+d^+_v) &
\log(1+d^-_v+d^+_v) &
1
\end{bmatrix}.
\]
Stacking all node features gives $X_p\in\mathbb{R}^{|V_p|\times 4}$.

Although the extracted call graph is directed, the current implementation symmetrizes it before message passing. If $A_p$ is the directed adjacency matrix, the model uses
\[
\tilde{A}_p = A_p + A_p^\top + I,
\]
where $I$ adds self-loops. It then applies symmetric normalization
\[
\hat{A}_p = \tilde{D}_p^{-1/2}\tilde{A}_p\tilde{D}_p^{-1/2},
\qquad
\tilde{D}_{p,ii}=\sum_j \tilde{A}_{p,ij}.
\]
This normalization is mathematically standard for GCNs, but it also removes caller-callee directionality, which later becomes an important limitation.

\subsection{GCN Build Failure Model}

The graph encoder uses two graph convolution layers:
\[
H_p^{(1)}=\operatorname{ReLU}(\hat{A}_pX_pW_0+b_0),
\]
\[
H_p^{(2)}=\operatorname{ReLU}(\hat{A}_pH_p^{(1)}W_1+b_1).
\]
The project embedding is the mean of final node states:
\[
g_p=\frac{1}{|V_p|}\sum_{v\in V_p}H_{p,v}^{(2)}.
\]
For a build $i$ belonging to project $p(i)$, let $z_i$ be its normalized CI metadata vector. The classifier predicts a logit
\[
s_i = f_\theta([g_{p(i)};z_i]),
\]
where $[\cdot;\cdot]$ denotes concatenation and $f_\theta$ is a two-layer multilayer perceptron. The predicted failure probability is
\[
\hat{y}_i=\sigma(s_i)=\frac{1}{1+\exp(-s_i)}.
\]
The target label is
\[
y_i=\mathbb{1}[\texttt{num\_fail\_class}_i>0].
\]

Because positive and negative builds are imbalanced, training uses weighted binary cross-entropy:
\[
\mathcal{L}
=-\frac{1}{N}\sum_{i=1}^{N}
\left[
\alpha y_i\log \hat{y}_i
+(1-y_i)\log(1-\hat{y}_i)
\right],
\qquad
\alpha=\frac{N_-}{N_+},
\]
where $N_+$ and $N_-$ are the numbers of positive and negative training builds. A build is classified as failing when $\hat{y}_i\ge \tau$, where $\tau$ is selected on the validation split to maximize F1 over thresholds in $[0.2,0.8]$.


\subsection{Important Consequence of the Formulation}

The formulation above implies that all builds from the same project share the same graph embedding $g_p$. Thus, for two builds $i$ and $j$ from the same project,
\[
p(i)=p(j)\quad\Longrightarrow\quad g_{p(i)}=g_{p(j)}.
\]
Any build-specific variation must therefore come only from $z_i$, the CI metadata vector. The graph cannot encode which files changed, which methods were executed, or which tests covered the changed code for a particular build. This mathematical property is central to the empirical failure analyzed in Section~\ref{sec:evaluation}: the graph representation is static and project-level, so it can easily become a soft project identifier rather than a build-specific dependency signal.

\section{Experimental Evaluation}
\label{sec:evaluation}

In this section, we evaluate the implemented static GCN baseline for build failure prediction. The goal is not to claim state-of-the-art performance, but to quantify what can and cannot be learned from whole-project static call graphs plus CI metadata.

\subsection{Research Questions}

Our evaluation is guided by three research questions. RQ1, predictive performance, asks whether a lightweight static GCN can predict failing builds better than simple baselines. RQ2, dataset sensitivity, asks how stable the results are across the available metadata datasets. RQ3, graph failure mode, asks whether the graph representation encodes useful build-specific dependency information or mostly behaves like a project identifier.

\subsection{Evaluation Protocol}

We evaluate three logical datasets: \texttt{dataset\_init.csv}, \texttt{dataset.csv}, and \texttt{dataset\_filter.csv}. The first file uses an older metadata schema, so we normalize it into the current training schema by mapping \texttt{ts\_duration}, \texttt{passcount}, and \texttt{failcount} into \texttt{test\_suite\_duration\_s}, \texttt{num\_pass\_class}, and \texttt{num\_fail\_class}; because it has no stage field, we assign \texttt{stage\_id=unknown}. For all datasets, the binary target is
\[
y_i=\mathbb{1}[\texttt{num\_fail\_class}_i>0].
\]

To approximate deployment in CI, build records are sorted by timestamp. The first 70\% are used for training, the next 15\% for validation, and the final 15\% for testing. The validation set is used to select the decision threshold $\tau$ that maximizes F1 over thresholds in $[0.2,0.8]$.

\subsection{Metrics and Baseline}

Let TP, TN, FP, and FN denote true positives, true negatives, false positives, and false negatives. We report
\[
\text{Accuracy}=\frac{\mathrm{TP}+\mathrm{TN}}{\mathrm{TP}+\mathrm{TN}+\mathrm{FP}+\mathrm{FN}},
\]
\[
\text{Precision}=\frac{\mathrm{TP}}{\mathrm{TP}+\mathrm{FP}},
\qquad
\text{Recall}=\frac{\mathrm{TP}}{\mathrm{TP}+\mathrm{FN}},
\]
\[
\text{F1}=\frac{2\cdot\text{Precision}\cdot\text{Recall}}{\text{Precision}+\text{Recall}}.
\]
These metrics are more informative than accuracy because the failure label is highly imbalanced and the positive rate changes over time.

We compare against an all-positive baseline that predicts every build as failing. If the positive rate is $\pi=N_+/N$, then this baseline has
\[
\text{Accuracy}=\pi,\qquad
\text{Precision}=\pi,\qquad
\text{Recall}=1,
\]
and therefore
\[
\text{F1}_{\mathrm{all+}}=\frac{2\pi}{1+\pi}.
\]
This baseline is intentionally simple, but it is mathematically important: when the positive rate is high, always predicting failure can achieve a strong F1 score.

\subsection{Overall Results}

\begin{table*}[ht]
    \centering
    \caption{Cross-dataset performance of the static GCN build failure predictor.}
    \label{tab:cross_dataset_results}
    \begin{tabular}{llrrrrrrrr}
        \toprule
        Dataset & Split & Rows & Pos. Rate & Best Epoch & $\tau$ & Acc. & Prec. & Rec. & F1 \\
        \midrule
        \texttt{dataset\_init} & Train & 2,334 & 0.891 & 40 & 0.20 & 0.919 & 0.973 & 0.934 & 0.953 \\
        \texttt{dataset\_init} & Validation & 500 & 0.434 & 40 & 0.20 & 0.732 & 0.799 & 0.512 & 0.624 \\
        \texttt{dataset\_init} & Test & 501 & 0.425 & 40 & 0.20 & 0.585 & 0.508 & 0.728 & 0.598 \\
        \midrule
        \texttt{dataset} & Train & 5,157 & 0.902 & 34 & 0.20 & 0.926 & 0.963 & 0.954 & 0.959 \\
        \texttt{dataset} & Validation & 1,105 & 0.702 & 34 & 0.20 & 0.837 & 0.906 & 0.857 & 0.881 \\
        \texttt{dataset} & Test & 1,106 & 0.389 & 34 & 0.20 & 0.741 & 0.961 & 0.347 & 0.509 \\
        \midrule
        \texttt{dataset\_filter} & Train & 5,157 & 0.902 & 34 & 0.20 & 0.926 & 0.963 & 0.954 & 0.959 \\
        \texttt{dataset\_filter} & Validation & 1,105 & 0.702 & 34 & 0.20 & 0.837 & 0.906 & 0.857 & 0.881 \\
        \texttt{dataset\_filter} & Test & 1,106 & 0.389 & 34 & 0.20 & 0.741 & 0.961 & 0.347 & 0.509 \\
        \bottomrule
    \end{tabular}
\end{table*}

Table~\ref{tab:cross_dataset_results} answers RQ1 and RQ2. The model fits the training data and performs well on validation for the current datasets, but it does not generalize reliably to the final temporal test period. On \texttt{dataset} and \texttt{dataset\_filter}, test precision is high at 0.961, but recall falls to 0.347. In other words, when the model predicts failure it is usually correct, but it misses most actual failures.

\begin{table}[ht]
    \centering
    \caption{Test F1 compared with the all-positive baseline.}
    \label{tab:baseline_results}
    \begin{tabular}{lrr}
        \toprule
        Dataset & All-positive F1 & GCN F1 \\
        \midrule
        \texttt{dataset\_init} & 0.597 & 0.598 \\
        \texttt{dataset} & 0.560 & 0.509 \\
        \texttt{dataset\_filter} & 0.560 & 0.509 \\
        \bottomrule
    \end{tabular}
\end{table}

Table~\ref{tab:baseline_results} shows that the GCN is effectively tied with the all-positive baseline on \texttt{dataset\_init} and worse than it on the two current datasets. This is a strong negative result: a graph-based method should at least outperform a classifier with no graph information.

\subsection{Temporal and Project Distribution Shift}

The temporal split creates a substantial distribution shift. In \texttt{dataset} and \texttt{dataset\_filter}, the training period is dominated by \texttt{kafka}: 4,589 of 5,157 training rows, with a 0.964 positive rate. The test period contains many more \texttt{james}, \texttt{hadoop}, and \texttt{hbase} rows, with much lower or more mixed positive rates. For example, \texttt{james} contributes 422 test rows but has only a 0.052 positive rate, while \texttt{kafka} contributes 137 test rows with a 0.993 positive rate.

This matters because the graph embedding $g_p$ is fixed per project. If training associates $g_{\texttt{kafka}}$ with frequent failures and associates other project embeddings with fewer observed failures, the classifier can appear strong during validation but fail when the project mix changes.

\subsection{Per-Project Test Behavior}

\begin{table*}[ht]
    \centering
    \caption{Per-project test metrics for \texttt{dataset} and \texttt{dataset\_filter}.}
    \label{tab:per_project_results}
    \begin{tabular}{lrrrrrrrrr}
        \toprule
        Project & Rows & Pos. Rate & TP & TN & FP & FN & Prec. & Rec. & F1 \\
        \midrule
        \texttt{hadoop} & 287 & 0.467 & 6 & 152 & 1 & 128 & 0.857 & 0.045 & 0.085 \\
        \texttt{hbase} & 257 & 0.537 & 6 & 118 & 1 & 132 & 0.857 & 0.043 & 0.083 \\
        \texttt{james} & 422 & 0.052 & 1 & 397 & 3 & 21 & 0.250 & 0.045 & 0.077 \\
        \texttt{kafka} & 137 & 0.993 & 136 & 0 & 1 & 0 & 0.993 & 1.000 & 0.996 \\
        \texttt{karaf} & 3 & 0.000 & 0 & 3 & 0 & 0 & 0.000 & 0.000 & 0.000 \\
        \bottomrule
    \end{tabular}
\end{table*}

Table~\ref{tab:per_project_results} answers RQ3. The model almost perfectly captures \texttt{kafka} failures but misses nearly all failures in \texttt{hadoop}, \texttt{hbase}, and \texttt{james}. This behavior is consistent with a project-prior model, not a build-specific dependency model.

\subsection{Graph Structure Diagnostics}

\begin{table*}[ht]
    \centering
    \caption{Static call graph statistics. The largest component is reported as a percentage of nodes.}
    \label{tab:graph_stats}
    \begin{tabular}{lrrrrrr}
        \toprule
        Project & Nodes & Edges & Avg. Degree & Weak Comp. & Largest Comp. & Max Degree \\
        \midrule
        \texttt{hadoop} & 87,016 & 253,694 & 5.83 & 974 & 96.9\% & 5,828 \\
        \texttt{hbase} & 42,609 & 131,416 & 6.17 & 470 & 96.2\% & 3,207 \\
        \texttt{jackrabbit-oak} & 34,245 & 94,337 & 5.51 & 426 & 96.1\% & 2,122 \\
        \texttt{kafka} & 29,757 & 87,855 & 5.90 & 577 & 95.0\% & 2,889 \\
        \texttt{james} & 23,723 & 69,061 & 5.82 & 337 & 95.9\% & 2,062 \\
        \texttt{karaf} & 7,870 & 21,794 & 5.54 & 166 & 94.3\% & 597 \\
        \bottomrule
    \end{tabular}
\end{table*}

The graph statistics in Table~\ref{tab:graph_stats} explain why the graph signal is weak. These are large sparse graphs with one giant weak component and many small components. Since the model mean-pools over all nodes, localized information is diluted:
\[
g_p=\frac{1}{|V_p|}\sum_{v\in V_p}H_{p,v}^{(2)}.
\]
If a failure-relevant region contains $k$ nodes in a graph with $|V_p|$ nodes, its direct contribution to the mean is only $k/|V_p|$. For a project such as \texttt{hadoop}, even a 100-node affected region contributes approximately
\[
\frac{100}{87016}\approx 0.00115
\]
of the unweighted mean. Without changed-node indicators or coverage-guided pooling, the model has little chance to preserve this local signal.

\subsection{Summary of Findings}

The static GCN fails for three mathematically concrete reasons. First, $g_p$ is invariant across builds from the same project, so the graph cannot represent build-specific changes. Second, symmetrization replaces $A_p$ with $A_p+A_p^\top+I$, discarding the direction of calls. Third, mean pooling over large graphs compresses tens of thousands of node embeddings into one vector, which washes out localized fault information. These limitations explain why the model performs like a project-prior classifier rather than a dependency-aware failure predictor.

\section{Relate Work}

% [proposal] list of all related work to the proposed idea. You should explain each related work briefly, in two or three sentences.

% [1] R. Cheng, S. Wang, R. Jabbarvand, and D. Marinov, "Revisiting Test-Case Prioritization on Long-Running Test Suites"

% According to the ISSTA '24 paper titled "Revisiting Test-Case Prioritization on Long-Running Test Suites," existing TCP often take over 6.5 hours to run test due to complex dependency issues in large-scale systems. This study introduces the LRTS dataset, which we utilize to demonstrate the efficiency improvement of GNN based approaches compared to baselines. In addition to these standard baselines, in this field, there is a recent shift toward integrating semantic code analysis with more complex methods to further push past performance limits.

% [2] S. R. Kongarana, A. R. Akepogu, and R. R. P, "XGRL-TCP: An Explainable Graph-Based Reinforcement Learning Framework for Test Case Prioritization in CI"

% [3] S. Sowmyadevi, A. Alphy, "Graph neural network-based mutation-aware regression test ordering using code dependency graphs and execution traces"

% For example, XGRL-TCP utilizes reinforcement learning with Graph Attention Networks to achieve APFD of 89.3\%. Another recent approach leverages mutation information to enhance fault detection reaching an average APFD of 88.9\%. While these models achieve high accuracy, they often come with significant computational costs. Our project aims to strike this balance by efficiently capturing structural dependencies within the LRTS environment using R-GCN.

% [4] Z. Feng et al., "CodeBERT: A Pre-Trained Model for Programming and Natural Languages"

% To capture the semantic meaning of code, CodeBERT [Feng et al., 2020] introduces a pre-trained model that processes programming languages like natural text. We will use this technique to generate vector representations of code changes, capturing semantic details which are invisible to history-based techniques.

Our research builds upon and contrasts with several key areas in regression testing and code representation. We position our empirical study in the context of the following related works.

Long-running test suites and history-based TCP. According to the ISSTA '24 paper by Cheng et al., existing Test Case Prioritization (TCP) methods often face severe bottlenecks in large-scale systems, taking over 6.5 hours to run tests due to complex dependencies. Their study introduced the LRTS dataset and the LatestFail+CC baseline, which relies on execution history. While LatestFail+CC is fast, it fundamentally struggles with predicting failures in newly introduced code dependencies, leaving a gap for structural approaches.

Complex graph and reinforcement learning models. To capture deeper code dependencies, there is a recent shift toward integrating more complex neural architectures. For instance, Kongarana et al. introduced XGRL-TCP, which utilizes reinforcement learning with Graph Attention Networks to achieve an APFD of 89.3\%. Similarly, Sowmyadevi and Alphy leveraged mutation information and execution traces with GNNs to reach a high average APFD of 88.9\%. While these models achieve high theoretical accuracy, they depend heavily on the flawless extraction of dynamic mutation or execution data and require massive computational overhead, making them difficult to scale in legacy LRTS environments.

Semantic code representation versus structural features. To capture the semantic meaning of source code, models like CodeBERT have been widely adopted to process programming languages like natural text. Initially, we considered using CodeBERT to generate dense vector representations of code changes. However, running a massive pre-trained transformer model to extract semantic embeddings for every commit introduces unacceptable inference overhead for a rapid CI gatekeeper. 

Our positioning. Unlike the aforementioned approaches that push for maximum accuracy using heavy dynamic traces or massive semantic models, our project takes a fundamentally different path. We aim to strike a practical balance by utilizing a lightweight Graph Convolutional Network (GCN) that operates \textit{exclusively} on static call graphs and temporal CI metadata. By doing so, we empirically demonstrate the performance lower-bound of a purely static approach and highlight the critical extraction bottlenecks that heavy dynamic models often overlook in real-world industrial settings.

\section{Threats to Validity}

Internal validity. The primary threat to internal validity lies in our architectural simplifications and data extraction bottlenecks. To handle massive project graphs, we employed a simple mean-pooling strategy to generate project embeddings. This approach likely caused oversmoothing, washing out localized fault signals and contributing to the low Recall. Furthermore, our inability to extract comprehensive JaCoCo dynamic coverage was dependent on our specific build environment and the inherent fragility of legacy dependency versions. A highly customized, heavily engineered CI setup might successfully bypass these extraction bottlenecks and yield a richer graph topology. Additionally, parsing static ASTs limits our visibility into dynamic dispatch and runtime polymorphism, which are common in object-oriented Java systems.

External validity. Our evaluation is restricted to the LRTS dataset, specifically 10 large-scale Apache open-source projects (e.g., Hadoop, Kafka). While these represent complex, real-world systems, their architecture and Java/Scala-based CI workflows might not perfectly generalize to proprietary, closed-source industrial repositories, microservice architectures, or projects written in other programming languages (e.g., C++ or Python).

Construct validity. A potential threat to construct validity arises from our pivot in the prediction target. By predicting whether a build contains at least one failing test class ($num\_fail\_class > 0$), we established a binary proxy for CI optimization. While this formulation effectively measures the model's capability to act as a preliminary gatekeeper, it does not fully capture the continuous time-saving metrics (such as APFD or APFDc) traditionally used in granular Test Case Prioritization (TCP) research.

\section{Conclusion}

In this project, we investigated the application of Graph Convolutional Networks (GCNs) for early build failure prediction in Long-Running Test Suites (LRTS). Initially aiming for a dynamic coverage-based approach, severe environment conflicts and data sparsity in legacy industrial systems forced a pivot to a lightweight GCN relying purely on static call graphs and temporal CI metadata. Our empirical evaluation demonstrated that while a static-only approach provides a highly scalable and computationally efficient preliminary gatekeeper, it suffers from significantly low Recall. This underwhelming performance highlights a critical semantic gap: static structures alone cannot adequately capture the complex runtime states that trigger dynamic regressions. Ultimately, our findings inadvertently validate the absolute necessity of dynamic analysis for accurate test prioritization. Future research must bridge this gap not merely by building heavier AI models, but by engineering more robust, low-overhead dynamic profiling tools that can seamlessly integrate into fragile, legacy industrial CI pipelines.
\end{document}
